\chapter{布尔电路}
\label{chap:boolean-circuits}

布尔电路（Boolean circuit）用有限个逻辑门（logic gate）把输入比特（bit）
变成输出比特。本讲
使用三种基本逻辑门：与门（AND）、或门（OR）和非门（NOT）。

\section{电路模型}

一个布尔电路 $C$ 可以看成有限的有向无环图（directed acyclic graph, DAG）。
图中包含：

\begin{enumerate}[label=\arabic*.]
  \item $n$ 个输入结点 $x_0,\ldots,x_{n-1}$；
  \item 若干内部结点，每个结点标记为 AND、OR 或 NOT 门；
  \item $m$ 个指定的输出结点 $y_0,\ldots,y_{m-1}$。
\end{enumerate}

AND 和 OR 门各有两个输入，NOT 门有一个输入。因为图中没有有向环，所以
各个门可以按拓扑顺序（topological order）依次求值。电路的\emph{规模}
（size）是其中逻辑门的数量，记为 $\card{C}$。

\begin{definition}[电路计算]
  \label{def:circuit-computes}
  设 $f:\bits^n\to\bits^m$。如果对每个输入
  $x\in\bits^n$，电路 $C$ 的 $m$ 个输出比特恰好组成 $f(x)$，则称
  $C$ \emph{计算}函数 $f$。
\end{definition}

\begin{example}
  \label{ex:xor-circuit}
  异或函数（exclusive OR, XOR）可以写成
  \[
    \operatorname{XOR}(a,b)
    =\neg(a\land b)\land(a\lor b),
  \]
  因而可以用 AND、OR 和 NOT 门组成的布尔电路计算。
\end{example}

\section{AON-Circ 程序与 NAND 门}

\begin{definition}[AON-Circ 程序]
  \label{def:aon-circ-program}
  AON-Circ 程序是一列按顺序执行的赋值语句（assignment statement）。
  每条语句调用一次 AND、OR
  或 NOT，并且只能读取输入变量或此前语句已经算出的变量。程序最后指定
  $m$ 个变量作为输出。程序的规模是赋值语句的行数。
\end{definition}

例如，一条语句可以具有以下三种形式之一：
\[
  z_i\gets z_j\land z_k,
  \qquad
  z_i\gets z_j\lor z_k,
  \qquad
  z_i\gets\neg z_j,
\]
其中右侧的变量必须已经定义。

\begin{theorem}
  \label{thm:circuit-program-equivalence}
  函数 $f$ 能由布尔电路计算，当且仅当它能由 AON-Circ 程序计算。两种表示
  可以保持规模不变：电路的门数等于对应程序的行数。
\end{theorem}

\begin{proof}
  给定电路，按拓扑顺序为每个门写一条赋值语句，即得到 AON-Circ 程序。
  反过来，为程序的每一行建立一个门，并把该行读取的变量连接到这个门，
  即得到布尔电路。每个门与每行语句一一对应。
\end{proof}

与非门（NAND）定义为
\[
  \operatorname{NAND}(a,b)=\neg(a\land b).
\]
只使用 NAND 门便可实现三个基本门：
\begin{align*}
  \neg a
    &=\operatorname{NAND}(a,a),\\
  a\land b
    &=\operatorname{NAND}\bigl(
        \operatorname{NAND}(a,b),
        \operatorname{NAND}(a,b)
      \bigr),\\
  a\lor b
    &=\operatorname{NAND}\bigl(
        \operatorname{NAND}(a,a),
        \operatorname{NAND}(b,b)
      \bigr).
\end{align*}
因此 NAND 电路与普通布尔电路可以相互转换。若原电路规模为 $s$，则转换后
分别满足
\[
  \card{C_{\mathrm{AON}}}\leq 2\card{C_{\mathrm{NAND}}},
  \qquad
  \card{C_{\mathrm{NAND}}}\leq 3\card{C_{\mathrm{AON}}}.
\]

\section{任意布尔函数的电路实现}

\begin{theorem}
  \label{thm:all-boolean-functions-have-circuits}
  对任意 $n,m\geq 1$ 以及任意函数
  \[
    f:\bits^n\longrightarrow\bits^m,
  \]
  都存在计算 $f$ 的布尔电路。进一步，电路规模可以达到
  \[
    \card{C}=O\!\left(\frac{m2^n}{n}\right).
  \]
\end{theorem}

\paragraph{粗略上界的构造.}
先考虑 $f$ 的第 $j$ 个输出比特 $f_j$。对每个
$a=(a_0,\ldots,a_{n-1})\in\bits^n$，定义
\[
  T_a(x)=\bigwedge_{r=0}^{n-1}\ell_{a,r}(x_r),
  \qquad
  \ell_{a,r}(x_r)=
  \begin{cases}
    x_r, & a_r=1,\\
    \neg x_r, & a_r=0.
  \end{cases}
\]
当且仅当 $x=a$ 时，$T_a(x)=1$。因此
\[
  f_j(x)=\bigvee_{\substack{a\in\bits^n\\f_j(a)=1}}T_a(x).
\]
这就是析取范式（disjunctive normal form, DNF）的直接构造。每个 $T_a$
使用 $O(n)$ 个门，而这样的项至多有 $2^n$ 个，所以一个输出
比特需要 $O(n2^n)$ 个门。分别构造 $m$ 个输出，便得到
\[
  \card{C}=O(mn2^n).
\]
\cref{thm:all-boolean-functions-have-circuits}中的
$O(m2^n/n)$ 上界去掉了这个直接构造中的冗余。这个更精细的上界将在后文
证明。

\begin{definition}[通用门组]
  \label{def:universal-gate-set}
  如果只使用门集合 $\mathcal F$ 中的门就能计算任意布尔函数，则称
  $\mathcal F$ 是\emph{通用的}（universal）。
\end{definition}

由\cref{thm:all-boolean-functions-have-circuits}和上面的三个恒等式，NAND
门本身就是通用的。因此，要判断一组门 $\mathcal F$ 是否通用，只需判断
能否仅用 $\mathcal F$ 中的门实现 NAND 函数。

为了证明\cref{thm:all-boolean-functions-have-circuits}，我们先来介绍一些
特性和例子。

\subsection{NAND-Circ 程序中的语法糖（syntactic sugar）}

为了让程序更容易书写，可以暂时加入以下三种\emph{语法糖}：

\begin{enumerate}[label=\arabic*.]
  \item \textbf{定长循环（fixed-length loop）：}把循环体复制固定次数；
  \item \textbf{用户定义过程（user-defined procedure）：}在调用处展开过程体；
  \item \textbf{条件语句（conditional statement）：}同时计算两个分支，
    再用选择器（multiplexer, MUX）输出其中一个。
\end{enumerate}

条件语句中的一比特选择器可以写成
\[
  \operatorname{MUX}(b,u,v)
  =(\neg b\land u)\lor(b\land v),
\]
其中 $b=0$ 时输出 $u$，$b=1$ 时输出 $v$。循环次数和过程参数一旦固定，
上述写法都能展开成普通的 NAND-Circ 程序，因此不会增强计算能力。

\subsection{ADD 与 LOOKUP}

\begin{example}[ADD]
  \label{ex:add-circuit}
  输入包含两个 $n$ 位二进制数。按从低位到高位的顺序，将它们分别记为
  \[
    (x_0,\ldots,x_{n-1}),
    \qquad
    (x_n,\ldots,x_{2n-1})
  \]
  \par\noindent
  加法函数为
  \[
    \operatorname{ADD}_n:\bits^{2n}\longrightarrow\bits^{n+1}.
  \]
  令 $\operatorname{MAJ}(a,b,c)$ 表示三个比特的多数值（majority），即
  \[
    \operatorname{MAJ}(a,b,c)
    =(a\land b)\lor(a\land c)\lor(b\land c).
  \]
  可以用如下定长循环实现逐位进位加法：
\begin{verbatim}
result = [0] * (n + 1)
carry  = [0] * (n + 1)
for i in range(n):
    result[i]  = XOR(carry[i], XOR(x[i], x[i+n]))
    carry[i+1] = MAJ(carry[i], x[i], x[i+n])
result[n] = carry[n]
return result
\end{verbatim}
  XOR 和 MAJ 都能用常数个 NAND 门实现。把循环展开后，程序共有 $O(n)$
  行，因而加法电路的规模为 $O(n)$。
\end{example}

\begin{example}[LOOKUP]
  \label{ex:lookup-circuit}
  输入由 $2^k$ 个数据比特
  $x_0,\ldots,x_{2^k-1}$ 和 $k$ 个索引比特
  $i_0,\ldots,i_{k-1}$ 组成。令
  \[
    j=\sum_{r=0}^{k-1}i_r2^{k-1-r}.
  \]
  查表函数（lookup function）返回第 $j$ 个数据比特：
  \[
    \operatorname{LOOKUP}_k:
    \bits^{2^k+k}\longrightarrow\bits,
    \qquad
    \operatorname{LOOKUP}_k(x,i)=x_j.
  \]
  它可以递归地把数据表分成两半：
\begin{verbatim}
lookup(x, index):
    if len(index) == 1:
        if index[0] == 0: return x[0]
        else:             return x[1]
    half = len(x) / 2
    if index[0] == 0:
        return lookup(x[0:half], index[1:])
    else:
        return lookup(x[half:], index[1:])
\end{verbatim}
  对电路而言，条件语句的两个分支都要构造，再由 MUX 选择输出。若
  $L(k)$ 表示展开后的程序行数，则
  \[
    L(k)=2L(k-1)+O(1),
    \qquad
    L(1)=O(1).
  \]
  因此 $L(k)=O(2^k)$，LOOKUP 电路的规模也是 $O(2^k)$。
\end{example}

\subsection{精细上界}

\begin{proof}[\cref{thm:all-boolean-functions-have-circuits}的精细上界证明]
  先考虑单输出函数 $f:\bits^n\to\bits$。取 $k<n$，并把输入写成

  \[
    x=(u,v),
    \qquad
    u\in\bits^k,
    \quad
    v\in\bits^{n-k}.
  \]

  对每个 $v\in\bits^{n-k}$，定义 $k$ 元函数

  \[
    h_v(u)=f(u,v).
  \]

  先共享地构造全部 $2^k$ 个最小项（minterm）

  \[
    M_a(u)=\bigwedge_{r=0}^{k-1}
    \begin{cases}
      u_r, & a_r=1,\\
      \neg u_r, & a_r=0,
    \end{cases}
    \qquad a\in\bits^k,
  \]

  代价为 $O(k2^k)$。任意函数 $h:\bits^k\to\bits$ 都可以写成

  \[
    h(u)=\bigvee_{\substack{a\in\bits^k\\h(a)=1}}M_a(u),
  \]

  因此，在最小项已经构造好的前提下，每个 $h$ 只需额外 $O(2^k)$ 个门。
  $k$ 元布尔函数总共只有 $2^{2^k}$ 种，所以全部不同的 $h_v$ 可以用

  \[
    O\!\left(k2^k+2^{2^k}2^k\right)
  \]

  个门计算。重复出现的 $h_v$ 直接复用同一根输出线。由此得到真值表
  （truth table）的一整行

  \[
    G(u)=\bigl(h_v(u)\bigr)_{v\in\bits^{n-k}}
    \in\bits^{2^{n-k}}.
  \]

  最后用 $v$ 查表：

  \[
    f(u,v)=\operatorname{LOOKUP}_{n-k}\bigl(G(u),v\bigr).
  \]

  由\cref{ex:lookup-circuit}，这一步需要 $O(2^{n-k})$ 个门。总规模为

  \begin{equation}
    \label{eq:refined-upper-bound-cost}
    O\!\left(k2^k+2^{2^k}2^k+2^{n-k}\right).
  \end{equation}

  对充分大的 $n$，令 $d=n-2\log_2 n$，并选择整数 $k$ 使得

  \[
    \frac d2<2^k\leq d.
  \]

  于是

  \[
    2^{2^k}2^k
      \leq \frac{2^n}{n^2}\,n
      =\frac{2^n}{n},
    \qquad
    2^{n-k}
      <\frac{2^{n+1}}{n-2\log_2 n}
      =O\!\left(\frac{2^n}{n}\right).
  \]

  同时 $k2^k=O(n\log n)=O(2^n/n)$，故
  \cref{eq:refined-upper-bound-cost}为 $O(2^n/n)$。
  有限多个较小的 $n$ 可以吸收到大 $O$ 的常数中。最后对 $m$ 个输出位
  分别使用这一构造，即得

  \[
    \card{C}=O\!\left(\frac{m2^n}{n}\right).
  \]
\end{proof}

\subsection{计数下界}

下面用计数论证（counting argument）说明上述关于 $n$ 的上界不能继续改进
到更低的数量级。

\begin{proposition}[布尔函数的电路下界]
  \label{prop:boolean-circuit-counting-lower-bound}
  存在常数 $c_0>0$，使得对所有充分大的 $n$，都有某个函数
  $f:\bits^n\to\bits$ 的任意布尔电路至少需要

  \[
    c_0\frac{2^n}{n}
  \]

  个逻辑门。
\end{proposition}

\begin{proof}
  一共有

  \[
    2^{2^n}
  \]

  个函数 $\bits^n\to\bits$。另一方面，设一个 NAND-Circ 程序至多有
  $s$ 行，并考虑 $s\geq n$ 的情形。程序中至多出现 $n+s\leq 2s$ 个值。
  即使把每行都宽松地记为一个三元组，并把所有不合法的三元组序列也计算
  在内，仍存在常数 $c\geq 1$，使这类程序的数量至多为

  \[
    2^{cs\log_2 s}.
  \]

  取

  \[
    s=\left\lfloor\frac{2^n}{2cn}\right\rfloor.
  \]

  当 $n$ 充分大时，$s\geq n$ 且 $\log_2s\leq n$，所以

  \[
    cs\log_2s\leq 2^{n-1},
    \qquad
    2^{cs\log_2s}<2^{2^n}.
  \]

  程序少于函数，故至少有一个函数不能由 $s$ 行以内的 NAND-Circ 程序
  计算。再由 NAND 门与 AND、OR、NOT 门之间的常数规模转换，结论成立。
\end{proof}

因此，\cref{thm:all-boolean-functions-have-circuits}给出的
$O(2^n/n)$ 单输出上界在常数因子意义下是紧的。

\section{程序编码与通用求值}

下面约定 $s\geq\max\{n,m,2\}$。一个含 $s$ 行、$n$ 个输入和 $m$ 个输出的
NAND-Circ 程序至多需要 $3s$ 个变量。把它们依次编号为

\[
  v_0,v_1,\ldots,v_{3s-1},
\]

其中 $v_0,\ldots,v_{n-1}$ 存放输入，$v_n,\ldots,v_{n+m-1}$ 存放输出。
每行程序编码为三元组

\[
  (a,b,c),
  \qquad
  v_a\gets\operatorname{NAND}(v_b,v_c).
\]

令 $\ell=\lceil\log_2(3s)\rceil$。每个下标用 $\ell$ 位表示，故整个程序
可以编码为长度 $3s\ell$ 的二进制串。

\begin{example}[通用求值函数]
  \label{ex:universal-evaluation}
  定义通用求值函数（universal evaluation function）

  \[
    \operatorname{EVAL}_{s,n,m}:
    \bits^{3s\ell+n}\longrightarrow\bits^m.
  \]

  输入写成 $(p,x)$，其中 $p$ 是一个程序编码，$x\in\bits^n$。规定

  \[
    \operatorname{EVAL}_{s,n,m}(p,x)=
    \begin{cases}
      P(x), & p\text{ 编码了合法程序 }P,\\
      0^m,  & \text{否则}.
    \end{cases}
  \]
\end{example}

\begin{theorem}[通用 NAND-Circ 程序]
  \label{thm:universal-nand-circ-program}
  对任意 $s\geq\max\{n,m,2\}$，都存在计算
  $\operatorname{EVAL}_{s,n,m}$ 的 NAND-Circ 程序 $U_{s,n,m}$，并且

  \[
    \card{U_{s,n,m}}=O(s^2\log s).
  \]
\end{theorem}

\begin{proof}
  以下出现的 LOOKUP、MUX 与等值测试都按前文展开为 NAND 门。
  用 $V\in\bits^{3s}$ 保存被模拟程序的全部变量值。首先令

  \[
    V^{(0)}=(x_0,\ldots,x_{n-1},0,\ldots,0).
  \]

  对一个用 $\ell$ 位表示的下标 $i$，将 $V$ 补零到长度 $2^\ell$，再调用
  \cref{ex:lookup-circuit}，即可实现

  \[
    \operatorname{GET}(V,i)=V_i
  \]

  ，代价为 $O(2^\ell)=O(s)$。更新操作逐坐标定义为

  \[
    \operatorname{UPDATE}(V,i,z)_j
    =\operatorname{MUX}\bigl(
      \operatorname{EQ}_{\ell}(i,j),V_j,z
    \bigr),
    \qquad 0\leq j<3s,
  \]

  其中 $\operatorname{EQ}_{\ell}$ 是 $\ell$ 位等值测试（equality test），
  可用 $O(\ell)=O(\log s)$ 个门实现。因此

  \[
    \card{\operatorname{UPDATE}}=O(s\log s).
  \]

  若程序的第 $t$ 行编码为 $(a_t,b_t,c_t)$，则依次计算

  \begin{align*}
    r_t&=\operatorname{GET}(V^{(t-1)},b_t),\\
    q_t&=\operatorname{GET}(V^{(t-1)},c_t),\\
    z_t&=\operatorname{NAND}(r_t,q_t),\\
    V^{(t)}&=\operatorname{UPDATE}(V^{(t-1)},a_t,z_t).
  \end{align*}

  每轮使用 $O(s\log s)$ 个门，模拟 $s$ 行共需 $O(s^2\log s)$ 个门。
  同时维护“变量是否已经定义”的比特，并检查下标范围、读取顺序和输出是否
  已定义，只会使每轮多出常数次 GET、UPDATE 与等值测试。最后用合法性比特
  屏蔽输出：合法时输出 $V_n^{(s)},\ldots,V_{n+m-1}^{(s)}$，否则输出
  $0^m$。总规模仍为 $O(s^2\log s)$。
\end{proof}

\begin{question}
  能否把\cref{thm:universal-nand-circ-program}的规模改进到
  $O(s\log s)$？
\end{question}
